% ====================== TÓM TẮT ======================
\chapter*{Tóm tắt}
\addcontentsline{toc}{chapter}{Tóm tắt}

Phân đoạn ngữ nghĩa (\textit{semantic segmentation}) ảnh chụp từ thiết bị bay không người lái (UAV) là bài toán nền tảng cho các ứng dụng giám sát và kiểm kê tài nguyên rừng. Tuy nhiên, dữ liệu rừng thu thập từ UAV gặp thách thức nghiêm trọng về \textbf{mất cân bằng lớp}: tỷ lệ chênh lệch điểm ảnh giữa lớp phổ biến nhất và lớp hiếm nhất lên tới hơn \textbf{2.700 lần} (thảm thực vật chiếm $35{,}4\%$ trong khi hàng rào chỉ chiếm $0{,}01\%$). Hệ quả là các mô hình huấn luyện bằng hàm mất mát Cross-Entropy truyền thống có xu hướng bỏ sót các lớp đối tượng hiếm nhưng quan trọng như hàng rào, ô tô hay cây đổ.

Báo cáo này nghiên cứu và đánh giá hệ thống hai chiến lược xử lý mất cân bằng lớp ở cấp độ hàm mất mát: (1) hàm mất mát Cross-Entropy với trọng số lớp ENet cố định làm chuẩn đối sánh (Exp 2); và (2) hàm mất mát kết hợp Focal + Dice + Cross-Entropy với trọng số thích nghi tự động học theo cơ chế độ bất định đa nhiệm (Exp 4). Các thực nghiệm được triển khai toàn diện trên bốn kiến trúc mạng phân đoạn phổ biến: DeepLabV3+, SegFormer, UNet và UNet++ trên bộ dữ liệu kết hợp thời tiết phức tạp \textit{sunny+overcast} (18 chuỗi ảnh, 26.242 cặp ảnh).

Kết quả thực nghiệm trên tập kiểm thử độc lập cho thấy chiến lược học thích nghi (Exp 4) cải thiện vượt trội điểm số mIoU từ $+4{,}15$ pp đến $+4{,}88$ pp trên toàn bộ các kiến trúc, trong đó mạng UNet++ đạt hiệu năng cao nhất với **mIoU $= 0{,}8374$** (so với $0{,}7956$ của Exp 2). Đặc biệt, hiệu năng trên các lớp hiếm tăng trưởng mạnh mẽ (hàng rào tăng từ $0{,}5581$ lên $0{,}7031$; ô tô tăng từ $0{,}7938$ lên $0{,}8715$; cây đổ tăng từ $0{,}5848$ lên $0{,}6707$). Báo cáo cũng chỉ ra sự đánh đổi quan trọng khi tối ưu hóa IoU vùng làm tăng nhẹ tỷ lệ bỏ sót (Miss Rate) đối với các vật thể có kích thước cực nhỏ. Cuối cùng, khả năng tổng quát hóa vượt trội của phương pháp đề xuất được kiểm chứng thành công trên bộ dữ liệu thực tế ngoài đời thực \textit{WildUAV}, vượt qua kết quả công bố của tác giả gốc.

\vspace{0.3cm}
\noindent\textbf{Từ khóa:} phân đoạn ngữ nghĩa, ảnh UAV, giám sát rừng, mất cân bằng lớp, Focal Loss, Dice Loss, học trọng số thích nghi, UNet++, WildUAV.

\newpage
% ====================== ABSTRACT ======================
\chapter*{Abstract}
\addcontentsline{toc}{chapter}{Abstract}

Semantic segmentation of aerial images captured by Unmanned Aerial Vehicles (UAVs) is a fundamental task for forest monitoring and inventory applications. However, UAV forest imagery exhibits a severe \textbf{class imbalance} challenge: the pixel ratio between the most dominant class and the rarest class exceeds \textbf{2,700 times} (ground vegetation at $35.4\%$ versus fence at $0.01\%$). Consequently, models trained with standard Cross-Entropy loss tend to ignore rare but critical classes such as fences, cars, and fallen trees.

This report presents a systematic evaluation of two class-imbalance mitigation strategies at the loss function level: (1) ENet-weighted Cross-Entropy loss as the baseline (Exp 2); and (2) a compound Focal + Dice + Cross-Entropy loss with adaptive weights automatically learned via multi-task uncertainty weighting (Exp 4). The experiments are comprehensively conducted across four representative segmentation architectures: DeepLabV3+, SegFormer, UNet, and UNet++ on a combined multi-weather dataset consisting of \textbf{18 sequences} (26,242 image pairs).

Experimental results on an independent test set demonstrate that the adaptive loss strategy (Exp 4) consistently outperforms the baseline, improving mIoU by $+4.15$ pp to $+4.88$ pp across all models. UNet++ achieves the highest performance with a **test mIoU of $0.8374$** (compared to $0.7956$ in Exp 2). Specifically, performance on rare classes increases significantly (fence IoU rises from $0.5581$ to $0.7031$; car IoU from $0.7938$ to $0.8715$; and fallen trees from $0.5848$ to $0.6707$). The report also highlights a crucial trade-off: optimizing for region overlapping (IoU) slightly increases the miss rate for extremely small objects. Finally, the superior generalization and robustness of the proposed method are validated on the real-world \textit{WildUAV} dataset, where our model successfully outperforms the original publication's baseline.

\vspace{0.3cm}
\noindent\textbf{Keywords:} semantic segmentation, UAV imagery, forest monitoring, class imbalance, Focal Loss, Dice Loss, uncertainty weighting, UNet++, WildUAV.
